\section{Consequences and limits}

For \(a=b=2\), the initial domain is
\[
\{(x,y)\in\D^2:2|xy|<1\},
\]
and the continued function has the polar components
\(2x^iy^j=1\) for \((i,j)\in\prim\).
This specialization illustrates both distinctions in the theorem:
the diagonal series is rational, but the joint bidisc boundary is
impenetrable; and the curved boundary of initial convergence is
not everywhere a barrier to continuation.
For example, \((x,y)=(\mathrm i/\sqrt2,\mathrm i/\sqrt2)\) lies on
\(2|xy|=1\) but on no polar component. The pair \((i,j)=(1,1)\)
gives \(2xy=-1\), while every other positive primitive pair has
\(i+j>2\) and hence \(|2x^iy^j|<1\).
The continued function is therefore holomorphic near this point.
This is a direct illustration of the classified divisor, not an
additional independent result.

The arithmetic and analytic inputs have separate roles.
Corvaja--Zannier supplies the subexponential two-index bound in the
independent branch. Elementary common-base algebra determines the
dependent convergence domain and ray exponents. Positive atomic
measures give one family of natural-boundary slices for both
branches. Classical subharmonic propagation turns those slices into
joint boundary obstructions where a holomorphic product is
available, and the actual dependent pole grids settle the remaining
parameter range.

The result does not assert that every fixed complex slice has a
natural boundary; the zero axes show why that would be false.
Nor does it give an effective exceptional threshold for the
\(S\)-unit theorem, a new diagonal recurrence classification, or a
new definition of rectangular returns.
It concerns the ordinary two-clock series \eqref{eq:intro-series}
and its native common fixed counts, without replacing the two
integer times by volume, subgroup index, or another arithmetic clock.
